\begin{pastexam}[label={px:ncu-107-dsa-12}]{中央資工所 107 學年度 第 12 題}
（單選；每題 5 分，答錯倒扣一分）

Program quickSort below aim to apply quick sort to A[left: right] into non-decreasing order.
\begin{minted}[fontsize=\small,breaklines]{c}
void quickSort ( int A[], int left, int right )
{
  int pivot, i, j;
  int temp;
  if ( left < right ) {
     i = left-1; j = right;
     pivot = A[right];
     do {
          do i++; while ( A[i] < pivot );
          do j--; while ( A[j] > pivot );
          If ( __(B4)__ ) swap(A[i], A[j], temp);   /* exchange A[i] and A[j] */
     } while ( __(B5)__ );
     swap ( __(B6)__ );
     quickSort( A, left, j-1 );
     quickSort( A, j+1, right ); }
}
\end{minted}

Blank (B5) in the algorithm above should be

(a)~\texttt{i>right-1} \qquad (b)~\texttt{j<left+1} \qquad (c)~\texttt{i>j} \qquad (d)~\texttt{i>=j} \qquad (e)~\texttt{i<j}
\end{pastexam}

\begin{pastsoln}
答案是 (e) \texttt{i<j}。

分割的過程見上一題（歷屆試題~\ref{px:ncu-107-dsa-11}）。 (B5) 是外層 \texttt{do}--\texttt{while} 繼續的
條件：$i$、$j$ 還沒交錯（$i < j$），表示中間還有沒看過的數，要再掃一回；一旦 $i \ge j$，左邊
$A[\mathit{left}..i-1]$ 都 $\le \mathit{pivot}$、$A[j+1..\mathit{right}-1]$ 都 $\ge \mathit{pivot}$，而
$i \ge j$ 使這兩段蓋住整個範圍，分割完成，離開迴圈。 上一題的追蹤：第 1 回合後 $i = 1 < j = 4$，繼續；
第 2 回合後 $i = 3 > j = 2$，停。

逐項檢查：
\begin{itemize}
\item (a) \texttt{i>right-1}：一開始 $i$ 很小，條件不成立，第 1 回合做完就離開，中間還沒掃完。 上面的例子
第 1 回合後就停（$i = 1$），接著把基準換到 $A[1]$，得 $4, 5, 2, 7, 9, 1$，$5$ 的右邊還有 $2$ 和 $1$，錯。
\item (b) \texttt{j<left+1}：同樣在第 1 回合後通常不成立，太早離開，錯。
\item (c) \texttt{i>j}、(d) \texttt{i>=j}：方向相反。 沒交錯時反而離開；交錯了還繼續掃，$i$、$j$ 會越過彼此
繼續走，甚至跑出範圍，錯。
\item (e) \texttt{i<j}：對。
\end{itemize}
\end{pastsoln}
